\mxnewpage\section*{Correction EV4-T1}\label{EV4-T1-CORR}
\textbf{QCM :} 1 B, 2 B, 3 B, 4 C, 5 C, 6 B, 7 C, 8 A, 9 A, 10 A, 11 A, 12 D. Chaque réponse correcte : 0{,}5 point ; total 6.
\paragraph{Vérifications des automatismes.} Priorités ; somme de fractions ; soustraction des exposants ; coefficient multiplicateur 1{,}2 ; division par 0{,}8 ; produit 1{,}25 fois 0{,}8 ; résolution affine ; accroissement vertical divisé par l’horizontal ; moyenne ; angle divisé par 360 ; complément à 1 ; substitution entre parenthèses.
\subsection*{Exercice 1 -- Choisir une loi}
\begin{enumerate}
\item $E(X)=2$ ; $E(Y)=-\dfrac{2}{3}+\dfrac{8}{3}=2$. \textit{Barème : 1 point. Résultat et justification élémentaire attendus ; demi-point possible pour démarche correcte inachevée.}
\item $E(Y^2)=\dfrac{2}{3}+\dfrac{64}{3}=22$, variance 18, contre 0 pour X. \textit{Barème : 2 points. Méthode ou justification : 1 ; calcul et conclusion : 1.}
\item Linéarité : espérance 4. X est en outre constant. \textit{Barème : 2 points. Méthode ou justification : 1 ; calcul et conclusion : 1.}
\end{enumerate}
\subsection*{Exercice 2 -- Seuil discret et modèle continu}
\begin{enumerate}
\item $Q\prime=-64e^{-t}<0$. \textit{Barème : 1 point. Résultat et justification élémentaire attendus ; demi-point possible pour démarche correcte inachevée.}
\item $Q(2)>8$, $Q(3)<8$ ; décroissance : premier rang 3. \textit{Barème : 2 points. Méthode ou justification : 1 ; calcul et conclusion : 1.}
\item Géométrique, premier 64, raison $e^{-1}$ ; somme $\dfrac{64(1-e^{-4})}{1-e^{-1}}$. \textit{Barème : 2 points. Méthode ou justification : 1 ; calcul et conclusion : 1.}
\end{enumerate}
\subsection*{Exercice 3 -- Équation et cercle}
\begin{enumerate}
\item $(x-1)^2+(y-2)^2=4$ ; centre (1;2), rayon 2. \textit{Barème : 1 point. Résultat et justification élémentaire attendus ; demi-point possible pour démarche correcte inachevée.}
\item (1;0) et (1;4). \textit{Barème : 1 point. Résultat et justification élémentaire attendus ; demi-point possible pour démarche correcte inachevée.}
\item Rayon vertical au point (1;4), tangente horizontale $y=4$. \textit{Barème : 2 points. Méthode ou justification : 1 ; calcul et conclusion : 1.}
\end{enumerate}
\paragraph{Étalonnage.} Durée cible : 20–25 min pour les automatismes, 80–85 min pour les exercices, 10–15 min de relecture. Ces durées sont des estimations éditoriales à mesurer en classe. Accepter toute méthode correcte ; ne pas sanctionner deux fois une erreur numérique propagée avec cohérence. Le barème proposé ne vaut pas décision d’harmonisation de l’équipe.
